% concatenated from maintext.tex, macros.tex, packages.tex
\documentclass[runningheads,orivec]{llncs}
    %% logics
\newcommand{\CS}{\textnormal{\textsf{CS}}}
\newcommand{\CSA}{\textnormal{$\mathsf{CS}\mathcal{A}$}}
\newcommand{\GL}{\textnormal{\textsf{GL}}} 
\newcommand{\CSM}{\textnormal{\textsf{CSM}}} 
\renewcommand{\S}{\textnormal{\textsf{S}}} 
\renewcommand{\P}{\mathsf{P}}
\newcommand{\ER}{\mathsf{ER}}
\newcommand{\IC}{\mathsf{ICP}}
\newcommand{\RP}{\mathsf{RP}}
\newcommand{\GLP}{\mathsf{GLP}}

    %% theories
\newcommand{\EA}{\mathsf{EA}}
\newcommand{\PA}{\mathsf{PA}}
\newcommand{\ZF}{\mathsf{ZF}}
\newcommand{\TA}{\mathsf{TA}}
\newcommand{\PRA}{\mathsf{PRA}}


    %% calculi
\newcommand{\labCS}{{\textnormal{labCS}^\infty}}
\newcommand{\labIC}{{\textnormal{labICP}^\infty}}
\newcommand{\labER}{{\textnormal{labER}^\infty}}



    %% aligning
\newcommand{\abs}{\vspace{4mm}}
\newcommand{\vsp}{\vspace{-6mm}}


    %% rule names
\newcommand{\Ax}{\textnormal{Ax}}
\newcommand{\wk}{\textnormal{wk}}
\newcommand{\subV}{\textnormal{sub}_{l}}
\newcommand{\cut}{(\textnormal{cut})}
\newcommand{\id}{(\textnormal{id})}
\newcommand{\idg}{(\textnormal{id}_g)}
\newcommand{\trans}{(\textnormal{trans})}
\newcommand{\transBi}{\textnormal{trans}_{\circ \bullet}}


    %% symbols in formulas
\newcommand{\Dmnd}{\Diamond}
\newcommand{\LLra}{\Longleftrightarrow}
\newcommand{\Lra}{\Leftrightarrow}
\newcommand{\Ra}{\Rightarrow}
\newcommand{\la}{\langle}
\newcommand{\ra}{\rangle}
\newcommand{\spa}{\;|\;}
\newcommand{\lra}{\leftrightarrow}
\newcommand{\grt}{\mathrel{\raisebox{.4ex}{\rotatebox[origin = c]{180}{$\triangle$}}\hspace{-1mm}}}

 
    %% misc
\newcommand{\R}{\mathcal{R}}
\newcommand{\Pw}{\mathcal{P}}
\newcommand{\Mw}{\mathcal{M}}


    %% formulas
\newcommand{\ph}{\varphi}
\newcommand{\A}{\mathcal{A}}


    %% functions
\newcommand{\Lab}{\mathfrak{L}}
\newcommand{\Nest}{\mathfrak{N}}
\newcommand{\PL}{\mathrm{PL}}


    %% sets and sequents
\newcommand{\La}{\mathcal{L}}
\newcommand{\WL}{\mathbb{WL}}
\newcommand{\Seq}{\mathfrak{S}}
\newcommand{\Br}{\mathfrak{B}}
\newcommand{\tree}{\mathcal{T}}
\newcommand{\LaBi}{\La^{\Box \triangle}}
\newcommand{\Prop}{\mathtt{Prop}}
\newcommand{\Var}{\mathtt{Lab}}
\newcommand{\Prov}{\textnormal{Prov}}










\let\proof\relax
\let\endproof\relax
\usepackage[leqno]{amsmath}

\usepackage{hyperref} 
 \usepackage{tikz-cd}
 \usepackage{tikz}
 \usetikzlibrary{positioning}
 \usepackage{amssymb}               
 \usepackage{graphicx}        
 \usepackage[T1]{fontenc}
 \usepackage{ebproof}   
 \usepackage{mdframed}
\usepackage{amsthm}
\usepackage{stmaryrd}
\usepackage{adjustbox}
\usepackage{soul}

\usepackage{xcolor}
\usetikzlibrary{arrows,shapes.geometric,shapes.symbols, shapes.arrows}

\usepackage{suetterl}

\usepackage{stackengine}
\usepackage{appendix} 

\usepackage{multicol}
 

 




\SetSymbolFont{stmry}{bold}{U}{stmry}{m}{n}

%Theorem environments
\makeatletter
\let\c@proposition\c@theorem
\let\c@corollary\c@theorem
\let\c@lemma\c@theorem
\let\c@definition\c@theorem
\let\c@example\c@theorem
\let\c@remark\c@theorem
\makeatother
%
\usepackage[T1]{fontenc}
% T1 fonts will be used to generate the final print and online PDFs,
% so please use T1 fonts in your manuscript whenever possible.
% Other font encondings may result in incorrect characters.
%
\usepackage{graphicx}
% Used for displaying a sample figure. If possible, figure files should
% be included in EPS format.
%
% If you use the hyperref package, please uncomment the following two lines
% to display URLs in blue roman font according to Springer's eBook style:
\usepackage{color}
\renewcommand\UrlFont{\color{blue}\rmfamily}
%\urlstyle{rm}
%
\begin{document} 
%
\title{A Non-Wellfounded and Labelled Sequent Calculus for Bimodal Provability Logic}
%
\titlerunning{A Non-Wellfounded and Labelled Sequent Calculus for Bimodal PRL}
% If the paper title is too long for the running head, you can set
% an abbreviated paper title here
%
\author{Justus Becker}
%
\authorrunning{J. Becker}
% First names are abbreviated in the running head.
% If there are more than two authors, 'et al.' is used.
%
\institute{University of Birmingham, Birmingham, UK \\
\email{jvb443@student.bham.ac.uk}}
%
\maketitle              % typeset the header of the contribution
%
\begin{abstract}
We present a labelled and non-wellfounded calculus for the bimodal provability logic $\CS$.
The system is obtained by modelling the Kripke-like semantics of this logic. As in~\cite{Das_et_al24}, we enforce the second-order property of converse wellfoundedness by using techniques from cyclic proof theory. 
We will prove soundness and completeness of this system with respect to the semantics and provide a primitive decision procedure together with a way to extract countermodels. 

\keywords{Provability logic \and Modal Logic \and Bimodal logic \and Labelled sequents \and Non-wellfounded proofs.}
\end{abstract}
%
%
\section{Introduction}
Modal provability logics originate from interpreting the $\Box$-modality as ``it is provable that'' in a sufficiently strong (usually arithmetic) theory. 
In \cite{Solovay76}, Robert Solovay proved completeness of the modal logic $\GL$ with respect to the provability of Peano Arithmetic ($\PA$) as well as the completeness of the modal logic~$\S$ (also referred to as $\mathsf{GLS}$) with respect to the true provability of $\PA$.
Since then, the concept of provability logic has been extended towards more and more complex logics: for instance, Japaridze's polymodal provability logic $\GLP$ \cite{Japaridze86} or logics of non-standard provability \cite{Feferman60,Henk16,Shavrukov94,Visser89}.
Bimodal provability logics extend unimodal logics, such as $\GL$, by also incorporating the provability of a second theory. 

A useful technique in the proof theory of modal logic is the usage of \emph{labelled sequents}, which are syntactically enriched sequents utilising the Kripke frame correspondence of modal logics (see \cite{Negri05,Simpson1994,Vigano2000}). 
Here, we use a two-sorted approach to labels by using two kinds of relational atoms (see \cite{GargGenoveseNegri_2012} for a multi-labelled system), which  allows us to capture two modalities simultaneously. As the labelled structure internalises part of the semantics, we can sometimes read out countermodels from the leaves of failed proof search, which is a common feature of labelled calculi.

Cyclic and more generally \emph{non-wellfounded proofs} were originally developed for the modal $\mu$-calculus \cite{Niwiński96} and later also adapted for first-order logic handling induction \cite{Brotherston05,BrotherstonSimpson_2010}. In the recent decade non-wellfounded proof theory has been also adapted for Arithmetic \cite{Simpson_2017} and other logics, such as propositional dynamic logic \cite{DochertyRowe_2019}.

In \cite{Shamkanov14}, Daniyar Shamkanov provides the first cyclic and non-wellfounded proof systems for $\GL$. Since then, a lot of variants of cyclic systems for similar logics, such as $\mathsf{iGL}$ \cite{Iemhoff16,SierraMiranda23} and $\mathsf{Grz}$ \cite{Savateev17}, have been developed. For bimodal provability logic, however, no sequent system has been introduced thus far. 

We will combine these notions of non-wellfounded proofs together with labelled sequents, similarly to \cite{Das_et_al24}, to obtain a proof system for the basic bimodal provability logic Carlson-Smoryński ($\CS$), named after Tim Carlson and Craig Smoryński due to their foundational work in the area \cite{Carlson86,Smorynski85}.
Like in \cite{Das_et_al24}, we utilise the progress condition of our proofs to implement the second-order condition for models of the logic (i.e.~converse wellfoundedness).


In the following section, we introduce the syntax and Kripke-like semantics for $\CS$. In section \ref{sec:labCS}, we introduce the rules of a non-wellfounded and labelled proof system for this logic, and prove semantic soundness and completeness. The latter is obtained by a countermodel construction, which we get by a primitive proof search strategy, giving us a decision procedure for $\CS$. 


\section{Bimodal Provability Logic}

We work in a two-sorted modal propositional language with modalities,~$\Box$ and~$\triangle$. As we are in a classical setting, it is enough to consider $\to$ and $\bot$ as the only logical connectives (where $\bot$ is a connective of arity zero), besides a countable set of propositional atoms $\Prop = \{p,q,r,... \}$. Wellformed formulas are then created by the following grammar:
$$A::= \bot \spa p \spa A \to A \spa \Box A \spa \triangle A$$ 
%
with $p \in \Prop$. We call this language $\LaBi$. 
%We use further connectives as abbreviations: $\neg A := A \to \bot$, $A \vee B := \neg A \to B$, $A \wedge B := \neg (\neg A \vee \neg B)$, $\Dmnd A : = \neg \Box \neg A$, and $ \grt A := \neg \triangle\neg A$.

Let us now define the \emph{basic bimodal provability logic} $\CS$. It is basic insofar as any bimodal provability logic contains $\CS$.
We refer the reader to \cite[Chapter 8]{Artemov2005} for an extensive overview of bimodal provability logics.
The logic $\CS$ is defined as follows.

\begin{definition} \label{def:CS}
    The logic $\mathsf{CS}$ is the smallest subset of $\LaBi$ containing all classical tautologies, the following axiom schemas \vspace{-1em}
    \begin{multicols}{2} 
   \begin{itemize} \setlength{\itemindent}{2em}
    \item[$(\textnormal{K}_\Box)$]
            $\Box (A \to B) \to (\Box A \to \Box B)$
    \item[$(\textnormal{L}_\Box)$] 
          $\Box (\Box A \to A ) \to \Box A$
    \item[$(\textnormal{4}_{\triangle\Box})$]
          $\Box A \to \triangle \Box A$ 
    \item[$(\textnormal{K}_\triangle)$] 
           $\triangle (A \to B) \to (\triangle A \to \triangle B)$
    \item[$(\textnormal{L}_\triangle)$] 
            $\triangle (\triangle A \to A ) \to \triangle A$ 
    \item[$(\textnormal{4}_{\Box\triangle})$]
        $\triangle A \to \Box \triangle A$ 
   \end{itemize}
  \end{multicols} \vspace{-1em}  
  \noindent
  and closed under modus ponens (mp) and both necessitation rules, (N$_\Box$) and (N$_\triangle$) (figure \ref{fig:HilbertRules}).
\end{definition}

\begin{figure}[t!]
\begin{mdframed}
    \centering 
    \begin{prooftree}
        \hypo{A \to B}
        \hypo{A}
        \infer[left label=(mp)]2{B}
    \end{prooftree} \hspace{23mm}
    \begin{prooftree}
        \hypo{A}
        \infer[left label=(N$_\Box$)]1{\Box A}
    \end{prooftree}  \hspace{23mm}
    \begin{prooftree}
        \hypo{A}
        \infer[left label=(N$_\triangle$)]1{\triangle A}
    \end{prooftree} 
\end{mdframed}
    \caption{Rules of the Hilbert Axiomatisation for $\CS$}
    \label{fig:HilbertRules}
\end{figure}

Observe that we can derive the axioms $\Box A \to \Box \Box A$ and $\triangle A \to \triangle \triangle A$ by the L{\"o}b and Kripke axioms ($L_\Box$, $K_\Box$ $L_\triangle$, and $K_\triangle$), which gives us essentially four different axioms that are reminiscent of the unimodal axiom $4$ $(\Box A \to \Box \Box A)$. 
Notice also that any unimodal fragment (either $\triangle$-free or $\Box$-free) of $\CS$ is exactly $\GL$ (up to renaming of modalities).

Recall that the logic $\GL$ is sound and complete with respect to the provability of $\PA$ \cite{Solovay76}. More formally, this means that we interpret modal formulas as arithmetic sentences such that the logical connectives are preserved and $\Box$ is interpreted as a provability predicate; then for any modal formula $A$ and any such interpretation $(\cdot)^*$ we have $\GL \vdash A$ iff $\PA \vdash A^*$. 
Similarly, we can describe bimodal provability logics in terms of the provability of two recursively enumerable (r.e.) theories $T,U$, where $\Box$ is interpreted as a provability predicate for $T$ and $\triangle$ as the provability predicate for $U$. 
It is shown in \cite[Theorem 3.15.ii]{Smorynski85} that there are theories $T,U$ extending Elementary Arithmetic such that their joint provability logic coincides with $\CS$. Notice, however, that these theories will not be natural, as the two modalities know essentially as little as possible about each other: the only interaction axioms are $\textnormal{4}_{\triangle\Box}$ and $\textnormal{4}_{\Box\triangle}$, which are purely consequences of the provable $\Sigma_1$-completeness of $T$ and $U$. 

Just as for $\GL$, there also exists for $\CS$ a well-behaved Kripke-like semantics. 
Recall that a model for $\GL$ (a \emph{$\GL$-model}) is a tuple $\la W ,{\prec}, V \ra$ with a non-empty set of points~$W$, a conversely wellfounded relation $\prec \; \subseteq W \times W$ and some valuation function $V : \Prop \to \Pw (W)$. We obtain a Carlson model by adding for each modality a subset into the structure.

\begin{definition}[\cite{Carlson86}]
    A \emph{Carlson model} $\Mw$ is a tuple $\la W, \prec , M_0 , M_1 , V \ra$ such that $\la W, \prec , V \ra$ is a $\mathsf{GL}$-model with two subsets $M_0 , M_1 \subseteq W$. 

    A modal formula $A \in \LaBi$ is evaluated on $\Mw$ inductively as follows.

    $\Mw , x \nVdash \bot$

    $\Mw , x \Vdash p$ iff $x \in V(p)$

    $\Mw , x \Vdash A \to B$ iff $\Mw , x \Vdash B$ or $\Mw , x \nVdash A$

    $\Mw , x \Vdash \Box A$ iff for all $y$ such that $x\prec y$ and $y \in M_0$ we have $\Mw , y \Vdash A$

    $\Mw , x \Vdash \triangle A$ iff for all $y$ such that $x\prec y$ and $y \in M_1$ we have $\Mw , y \Vdash A$
\end{definition}

If a formula $A$ holds at all points of a model $\Mw$, we write $\Mw \Vdash A$, and $\Mw \nVdash A$ if it does not. 
In the literature of bimodal provability logic, one also considers Carlson models which consist of either two kinds of relations (see for example \cite{Visser95}) or only one subset (see for example \cite{Beklemishev94}). This works especially well for extensions of $\CS$, however the Carlson models we introduced are more appropriate for considering the logic $\CS$ itself.
The syntax and semantics is then connected by the following completeness theorem.

\begin{theorem}[\cite{Smorynski85}]\label{thm:soundComplCS}
    For any formula $A \in \La^{\Box \triangle}$: $\mathsf{CS} \vdash A$ iff for all Carlson models $\Mw$: $\Mw \Vdash A$.
\end{theorem} 
\begin{proof}
    See \cite[Chapter 4, Theorem 3.5]{Smorynski85}, where $\CS$ can be defined as the $\Box$-free fragment of $\mathsf{PRL}_2$, that is, $\GL$ extended by two modalities which are weaker than $\Box$ and generally independent of each other.
\end{proof}


\section{A Bi-Labelled Sequent Calculus} \label{sec:labCS}

In this section, we present our bi-labelled non-wellfounded calculus $\labCS$, and show that it is sound and complete with respect to the logic $\CS$.

\subsection{The Calculus $\labCS$}

The calculus we present is based on the relational semantics introduced in the previous section, using labels. Labelled systems can be traced back to Kanger \cite{Kanger57}, who used them to introduce the first cut-free system for the modal logic $\mathsf{S5}$. However, only much later have labelled systems been applied widely to non-classical logics \cite{Gabbay_1996,Negri05,Simpson1994,Vigano2000}. 

The second proof theoretic tool we use are \emph{non-wellfounded proofs}. Those are a generalisation of cyclic proofs which have been introduced for the modal $\mu$-calculus \cite{Niwiński96} and later also adapted for other logics, such as propositional dynamic logic \cite{DochertyRowe_2019} or logics with inductive definitions \cite{Brotherston05,Simpson_2017}. 
While wellfounded proofs are often seen as top-down inductively definable trees, non-wellfounded proofs are defined as coinductive trees, allowing for infinitary branches. 
Our system is inspired by the labelled, non-wellfounded systems for $\GL$ and $\mathsf{IGL}$ (an intuitionistic version of $\GL$ with explicit $\Dmnd$) from \cite{Das_et_al24}.

For the language of our calculus, we use a countable set of symbols, also called \textit{labels} $\Var=\{x,y,z,...\}$ and two additional symbols $R$ and $S$. A \textit{labelled formula} is then formed by $x:A$ with $x \in \Var$ and $A \in \LaBi$, and a relational labelled formula (often called \textit{relational atom}) is of the form $xRy$ or $xSy$ with $x,y \in \Var$. A \textit{labelled sequent} $\R , \Gamma \Ra \Omega$ consists of a finite multiset of relational atoms $\R$ and two finite multisets of simple labelled formulas $\Gamma $ and~$\Omega$. As usual, we use a comma to indicate multiset union, meaning that $\Gamma, \Omega$ is the (disjoint) union of the multisets $\Gamma$ and $ \Omega$. 
We will also write $\Gamma , x :A$ (or $\R, xRy$) to mean the union of $\Gamma$ with the singleton multiset $\{x:A\}$ (respectively the union of $\R$ and $\{xRy\}$).
We write $\Seq$ to vary over labelled sequents, and we refer to the set of labels occurring in a sequent $\Seq$ as $Lab(\Seq)$.  

\begin{figure}[t]
    \begin{mdframed}
    \begin{center} \vspace{3mm}
        \begin{prooftree}
            \hypo{}
            \infer[left  label=Id ]1[]{\R , \Gamma , x : p \Ra x : p , \Omega}
        \end{prooftree} \hspace{6mm}
        \begin{prooftree}
            \hypo{}
            \infer[left  label=$\bot$ ]1[]{\R , \Gamma , x : \bot \Ra \Omega}
        \end{prooftree} \\[4mm]
        \begin{prooftree}
            \hypo{\R , \Gamma, x: A \Ra x : B , \Omega}
            \infer[left  label=$\to$R ]1[]{\R , \Gamma \Ra x : A \to B , \Omega}
        \end{prooftree} \hspace{6mm}
        \begin{prooftree}
            \hypo{\R , \Gamma \Ra x : A, \Omega}
            \hypo{\R , \Gamma, x: B \Ra \Omega}
            \infer[left  label=$\to$L ]2[]{\R , \Gamma, x: A \to B \Ra \Omega}
        \end{prooftree} \\[4mm]
        \begin{prooftree}
            \hypo{\R, x R y , \Gamma \Ra y : A , \Omega}
            \infer[left  label=$\Box$R ]1[$y!$ ]{\R, \Gamma \Ra x : \Box A , \Omega}
        \end{prooftree} \hspace{6mm}
        \begin{prooftree}
            \hypo{\R, x R y , x: \Box A , y : A, \Gamma \Ra \Omega}
            \infer[left  label=$\Box$L ]1[]{\R, x R y , x : \Box A , \Gamma \Ra \Omega}
        \end{prooftree} \\[4mm]
        \begin{prooftree}
            \hypo{\R, x S y , \Gamma \Ra y : A , \Omega}
            \infer[left  label=$\triangle$R]1[$y!$]{\R, \Gamma \Ra x : \triangle A , \Omega}
        \end{prooftree} \hspace{6mm}
        \begin{prooftree}
            \hypo{\R, x S y , x: \triangle A , y : A , \Gamma \Ra \Omega}
            \infer[left  label=$\triangle$L ]1[]{\R, x S y , x : \triangle A ,  \Gamma \Ra \Omega}
        \end{prooftree} \\[4mm]
        \begin{prooftree}
            \hypo{\R, x \circ y , y \bullet z , x  \bullet z, \Gamma \Ra \Omega}
            \infer[left  label=$\transBi$ ]1[($\circ , \bullet \in \{R,S\}$)]{\R, x \circ y , y \bullet z , \Gamma \Ra \Omega}
        \end{prooftree}
    \end{center} 
    \end{mdframed}
    \caption{The rules of the system $\labCS$}
    \label{fig:RulesLabCS}
\end{figure}

The rules of the system $\labCS$ are given in figure \ref{fig:RulesLabCS}. Intuitively, an atom $xRy$ corresponds semantically to a relation $x \prec y$ in a Carlson model such that $y \in M_0$, while $xSy$ corresponds to the same where $y \in M_1$ (see also definition \ref{def:SequInterpret}); this is similar to the semantics introduced in \cite{Visser95}. 
Notice that the rule $\transBi$ constitutes actually four different rules, all together corresponding to the transitivity of $\prec$. 
One could also define a labelled system for $\CS$ by considering labelled formulas of the form $x \prec y$, $x \in M_0$, and $x \in M_1$ instead of $xRy$ and $xSy$. While this would reduce the amount of rules we would have to deal with, we would add a lot of structure which is more difficult to handle when controlling proofs; for example, in defining proof search strategies.
The rules $\Box$R and $\triangle$R have the additional constraint, signalled by $y!$, that the new label $y$ (read bottom up) should be \emph{fresh}. This means that $y$ is not allowed to occur in the conclusion of the rule.  
The sequent $\R , \Gamma \Ra \Omega$, which always occurs as a subsequent in all of the rules, is called the \emph{context} of that rule, whereas the other formulas in the sequents are either called \emph{principal} (if they are in the conclusion of the rule) or \emph{auxiliary} (if they occur in one of the premisses).
Besides the rules themselves, it is a crucial feature to allow for infinitely long branches inside a derivation, which are constrained by a progress condition. 

\begin{definition}[Pre-Proof, Trace, Progress, Proof]
     A \emph{pre-proof} is a possibly infinite tree generated from rules of $\labCS$.

    A \emph{trace} along a branch $(\Seq_i)_{i<\omega}$ (with $\Seq_i = \R_i , \Gamma_i \Ra \Omega_i$) of a pre-proof is a sequence of labels $(x_i)_{i < \omega}$ s.t.~for all $i<\omega$: either 1.~$x_{i+1} = x_i$; or 2.~${x_i R x_{i+1} \in \R_i}$; or 3.~$x_i S x_{i+1} \in \R_i$. 
    
    An infinite trace is \emph{progressing} if it is not eventually constant (i.e.~ conditions 2.~and 3.~above apply together infinitely often). A branch $(\Seq_i)_{i<\omega}$ in a pre-proof is progressing if it has a progressing trace.
    
    A \emph{proof} in the system $\labCS$ is then given by a pre-proof where each branch is either finite or progressing. We write $\labCS \vdash \Seq$ if there is a proof with $\Seq$ at the root. A formula $A$ is derivable in $\labCS$, written ${\labCS \vdash A}$, if $\labCS \vdash \Ra x : A$ (for any $x \in \Var$). 
\end{definition}


\begin{figure}[t]
    \centering 
    \begin{adjustbox}{max width=\textwidth}
        \begin{prooftree}
        \hypo{\vdots}
        \infer1{\R, zRu, xRu, x: \Box ( \Box p \to p ) \Ra z : p, u : p}
        \infer[left label=trans$_{R R}$]1{\R, zRu, x: \Box ( \Box p \to p ) \Ra z : p, u : p}
        \infer[left label=$\Box$R]1{\R, x: \Box ( \Box p \to p ) \Ra z : p, z : \Box p}
        \hypo{}
        \infer[left label=Id]1{\R, x: \Box ( \Box p \to p ), z : p \Ra z : p}
        \infer[left label=$\to$L]2{xSy,yRz, xRz, x: \Box ( \Box p \to p ), z : \Box p \to p \Ra z : p}
        \infer[left label=$\Box$L]1{xSy,yRz, xRz, x: \Box ( \Box p \to p ) \Ra z : p}
        \infer[left label=trans$_{S R}$]1{xSy,yRz, x: \Box ( \Box p \to p ) \Ra z : p}
        \infer[left label=$\Box$R]1{xSy, x: \Box ( \Box p \to p ) \Ra y : \Box p}
        \infer[left label=$\triangle$R]1{x: \Box ( \Box p \to p ) \Ra x : \triangle \Box p}
        \infer[left label=$\to$R]1{\Ra x : \Box ( \Box p \to p ) \to \triangle \Box p}
    \end{prooftree}
    \end{adjustbox}
    \caption{A proof in $\labCS$ with $\R = xSy,yRz, xRz$.}
    \label{fig:CSproof}
\end{figure}

\begin{example} \label{ex:CSproof}
    Consider as an example a proof in $\labCS$, given in figure \ref{fig:CSproof}. Observe that the leftmost branch of this proof tree has to be infinite. However, we can find a progressing trace, starting with $x,y,z,u,...$ (leaving out the constant steps).
\end{example}

One way to read the progress condition is to consider countermodel extractions. A desired feature of a calculus is the ability to read out countermodels from a failed proof. This is made even easier by incorporating labels. A countermodels can usually be read out from a single branch that is failed. This is either the case by that branch getting stuck (i.e.~no rules can be applied any more) or by that branch continuing indefinitely. Considering example \ref{ex:CSproof}, we could read out a countermodel from the leftmost branch. By the progress condition however, this model will not be a model of $\CS$. therefore, we want to include such branches for our system. This will be made more explicit in the following subsections (theorems \ref{thm:soundCS} and \ref{thm:complCS}).

In non-wellfounded proof theory, there is usually a special focus on \emph{regular} proofs, which can be written as finite graphs possibly containing cycles. 
In this work, regular proofs do not play a role as we do not consider effectiveness. Instead, we focus on logical and proof-theoretic aspects, such as showing soundness and completeness.


\subsection{Soundness}

We will prove both soundness and completeness with respect to Carlson semantics. To do so, we generalise the semantic truth of formulas to sequents.

\begin{definition}[Sequent interpretation] \label{def:SequInterpret}
    Let $\Seq = \R , \Gamma \Ra \Omega$ be a sequent and $\Mw = \la W , \prec, M_0 , M_1 , V \ra$ a Carlson model. 
    Let $I : Lab(\Seq) \to W$ be a function from labels of the sequent to points in the model.
    We call $I$ a \emph{sequent interpretation} of $\Seq$ on $\Mw$ if the following condition is fulfilled: If $xRy \in \R$ then $I(x) \prec I(y)$ and $I(y) \in M_0$, and if $xSy \in \R$ then $I(x) \prec I(y)$ and $I(y) \in M_1$.
\end{definition}

Intuitively, a sequent interpretation maps every label to a point in a Carlson model such that the meaning of the relational atoms match their interpretation in the Carlson model. 
From here on, we refer to sequent interpretations as interpretations, which we use to determine whether a sequent holds in a model.

\begin{definition}[Semantic Truth for Sequents]
     Let $\Mw = \la W , \prec, M_0 , M_1 , V \ra$ be a model. We say that a sequent $\Seq = \R , \Gamma \Ra \Omega$ is true on $\Mw$ (written $\Mw \Vdash \Seq$) if the following holds for all interpretations $I $ of $\Seq$ on $\Mw$:
     If $\Mw , I(x) \Vdash A$ for all $x : A \in \Gamma$, then there is some $y : B \in \Omega$ with $\Mw, I(y) \Vdash B$.
\end{definition}

Similarly to the validity of formulas, we write $\CS \Vdash \Seq$ if the sequent $\Seq$ is valid on all Carlson models, and $\CS \nVdash \Seq$ if it is not. 
We now proceed to prove soundness of $\labCS$ using the following crucial local soundness lemma. This lemma essentially tells us that truth is always preserved in a rule application, both upwards and downwards. Note that downwards preservation of truth always holds for sound sequent systems.

\begin{lemma}[Local Soundness] \label{lem:FalsityCS}
    Let $\Mw$ be a Carlson model and let
    \[
    \frac{\Seq_1 ... \Seq_n}{\Seq}
    \]
    be a rule of $\labCS$ with at least one premiss, then: $\Mw \Vdash \Seq$ iff $\Mw \Vdash \Seq_i$ for all $1 \leq i \leq n$.
\end{lemma}
\begin{proof}
    Generally, the proof is a simple inspection on the structure of the rules and the semantic meaning of the connectives as well as the conditions on Carlson models. We prove by contraposition and consider the following key cases.

    $\Box$R: Let $\Mw = \la W , \prec , M_0 , M_1 , V \ra$ be a model with  ${\Mw \nVdash \R , \Gamma \Ra x : \Box A , \Omega}$. 
    We can then find an interpretation $I$, for which the following holds: $\Mw , I(z) \nVdash B$ for all $z : B \in \Omega \cup \{x : \Box A\}$, and $\Mw , I(z) \Vdash B$ for all $z : B \in \Gamma$. 
    By $\Mw , I(x) \nVdash \Box A$, we must find some point $\tilde{y} \in M_0 \subseteq W$ with $I(x) \prec \tilde{y}$ and $\Mw , \tilde{y} \nVdash A$. 
    Extend $I$ to $I^\prime$ by adding some fresh label $y$ to its domain with $I(y)=\tilde{y}$. Then, $I^\prime$ is an interpretation of $\R , xR y , \Gamma \Ra y : A , \Omega$ over $\Mw$ with: $\Mw , I(z) \nVdash B$ for all $z : B \in \Omega \cup \{y : A\}$, and $\Mw , I(z) \Vdash B$ for all $z : B \in \Gamma$. Therefore $\Mw \nVdash \R , xR y , \Gamma \Ra y : A , \Omega$.\\ 
    See that the same argument works in the other direction as well:
    By assuming $\Mw \nVdash \R , xR y , \Gamma \Ra y : A , \Omega$ there must be an interpretation $I^\prime$ such that $\Mw , I(z) \nVdash B$ for all $z : B \in \Omega \cup \{y : A\}$, and $\Mw , I(z) \Vdash B$ for all $z : B \in \Gamma$. By $I^\prime$ being an interpretation and $x R y$ we have $I^\prime (x) \prec I^\prime (y)$, therefore $\Mw , I^\prime (x) \nVdash \Box A$. By reducing $I^\prime$ to $I$, we have an interpretation for $\R , \Gamma \Ra x : \Box A , \Omega$ as $y$ was introduced freshly. Thus, $\Mw \nVdash \R , \Gamma \Ra x : \Box A , \Omega$.

    \sloppy
    $\triangle$L: Consider some model $\Mw = \la W , \prec , M_0 , M_1, V \ra$ for which we have ${\Mw \nVdash \R , xS y , \Gamma, x : \triangle A \Ra \Omega}$. Then, we can find an interpretation $I$ with the following properties: $\Mw , I(z) \nVdash B$ for all ${z : B \in \Omega}$, and ${\Mw , I(z) \Vdash B}$ for all ${z : B \in \Gamma \cup \{ x : \triangle A \}}$. By $I$ being an interpretation and $xS y$ being in the relation set of the sequent: $I(x) \prec I(y)$ and $I(y) \in M_1$. Therefore, by $\Mw , I(x) \Vdash \triangle A$ we get $ \Mw , I(y) \Vdash A$. So, $\Mw , I(z) \Vdash B$ for all $z : B \in \Gamma \cup \{x : \triangle A , y : A \}$, which makes $I$ an interpretation witnessing that $\Mw \nVdash \R , xS y , \Gamma ,x : \triangle A , y : A \Ra \Omega$. \\
    The other direction follows directly by observing that a model $\Mw$ with an interpretation $I$ which falsifies $\R , xS y , \Gamma, x : \triangle A , y :A \Ra \Omega$ must also falsify the strictly smaller sequent $\R , xS y , \Gamma, x : \triangle A \Ra \Omega$.

    \sloppy
    $\transBi$: Consider some model $\Mw = \la W , \prec , M_0 , M_1 , V \ra$ where ${\Mw \nVdash \R , x \circ y , y \bullet z , \Gamma \Ra \Omega}$ with $\circ , \bullet \in \{R,S\}$. So, there is an interpretation $I$ with $\Mw , I(z) \nVdash B$ for all $z : B \in \Omega$, and $\Mw , I(z) \Vdash B$ for all $z : B \in \Gamma$. By $I$ being an interpretation and $x \circ y , y \bullet z$ being in the relation set of the sequent, we have $I(x) \prec I(y) \prec I(z)$, with either $I(z) \in M_0$ if $\bullet = R$, or $I(z) \in M_1$ if $\bullet = S$. So, by transitivity of ${\prec}:$ $I(x) \prec I(z)$. Thus, $I$ is also an interpretation of $\R , x \circ y , y \bullet z , x \bullet z ,\Gamma \Ra \Omega$, and it therefore witnesses that $\Mw \nVdash \R ,  x \circ y , y \bullet z , x \bullet z ,\Gamma \Ra \Omega$. \\
    The other direction follows immediately just like in the previous case.
\end{proof}

The following observation follows directly from the proof and will be helpful for our soundness argument.

\begin{corollary} \label{cor:interpretation}
    Given an instance of a rule of $\labCS$, with conclusion $\Seq$ and premisses $\Seq_1 , ... , \Seq_n$: if an interpretation $I$ witnesses that $\Mw \nVdash \Seq$, then there is ab interpretation $I^\prime$ witnessing $\Mw \nVdash \Seq_i$ which is either equal to $I$ or extends $I$.
\end{corollary}

Notice that the proof does not use the fact that our models are conversely wellfounded. This will be used to prove the next theorem when considering the infinite branches of a proof, which also have not played a role yet.


\begin{theorem}[Soundness of $\labCS$] \label{thm:soundCS}
    For any sequent $\Seq$: If $\labCS \vdash \Seq$ then $\CS \Vdash \Seq$.
\end{theorem}
\begin{proof}
    Assume $\labCS \vdash \Seq$ and suppose for a contradiction that $\CS \nVdash \Seq$.
    Then, there is some Carlson model $\Mw = \la W , \prec , M_0 , M_1 , V \ra$ that falsifies $\Seq$, hence $\Mw \nVdash \Seq$. 
    As $\labCS \vdash \Seq$, it must be witnessed by some proof $\tree$, where every branch of $\tree$ is either finite or progressing. We now extract an infinite branch $(\Seq_i)_{i < \omega }$ from $\tree$ which will not be progressing, contradicting the fact that $\tree$ is a proof. We define the branch inductively by $\Seq_0 := \Seq$, and $\Seq_{i+1}$ is a premiss of $\Seq_i$ in $\tree$ such that $\Mw \nVdash \Seq_{i+1}$. By $\Mw \nVdash \Seq$ and local soundness (lemma \ref{lem:FalsityCS}) this branch is well defined. 
    
    Note that $(\Seq_i)_{i < \omega }$ cannot be finite as it would have to end with an initial sequent (either $\R , \Gamma , x : p \Ra x : p , \Omega$ or $\R , \Gamma , x : \bot \Ra \Omega$) as $\tree$ is a proof and an initial sequent cannot be falsified by a model. Thus, $(\Seq_i)_{i < \omega }$ is an infinite branch, which must therefore be progressing (again, because $\tree$ is assumed to be valid). 

    Call $I_i$ the sequent interpretation of $\Seq_i$ which is guaranteed by $\Mw \nVdash \Seq_i$. Note that $I_{i+1}$ can be obtained by possibly just extending $I_{i}$ (corollary \ref{cor:interpretation}), which allows us to define the limit interpretation $I$ for all $\Seq_i$ over $\Mw$.
    
    We show that $\prec$ cannot be conversely wellfounded.
    Consider some label $x_n$ occurring in the progressing trace  $(x_i)_{i < \omega }$ of $(\Seq_i)_{i < \omega }$ together with its corresponding point in the model $I_n(x_n)$. By the trace being progressive we can find some $m > n$ s.t.~ $x_n \circ x_m$ (for $\circ \in \{R,S\}$), thus $I_n(x_n) \prec I_m (x_m)$. 
    Further, this shows that for any label $x$ occurring in the trace, there is some distinct label $x^\prime$ also occurring in the trace such that $I(x) \prec I(x^\prime)$ as all $I_i$ are contained in $I$. 

    We can then find a chain of points $I(x) \prec I(x^\prime) \prec I(x^{\prime\prime}) \prec ...$ with $x=x_1$. Either this chain will have infinitely many distinct elements or it must eventually repeat an element. If it has infinitely many elements, that means $\prec$ has an infinite non-cyclic path. If elements do repeat, there must be some loop in $\prec$.
    In both cases, $\prec$ is not conversely wellfounded, meaning that $\Mw$ is not a Carlson model, a contradiction.
\end{proof}

\subsection{Completeness}

We continue with the proof of completeness. There are generally two ways to establish completeness for a sequent-style calculus. One can either show that the system is complete wrt.~the Hilbert axiomatisation by introducing the cut rule -- which is then shown to be admissible -- or wrt.~semantics, by considering failed proofs and showing that one can read out a countermodel from them. 
We will do the latter here.
As we have a fully invertible calculus, we are able to consider a technique which is relatively simple and has been inspired by an unpublished (personally communicated) completeness proof for the labelled and non-wellfounded system for classical $\GL$ from \cite{Das_et_al24}.


Let us define a proof search procedure which we define in two phases. Dividing the procedure like this makes it easy to show that the strategy is well defined. It also allows us to adapt the strategy to other logics, for example if we have a relational rule that can fire infinitely often such as directedness for the modal logic $\mathsf{KD}$. We divide the phases by considering firstly rules that do not introduce new labels, which we will call \emph{saturation}; secondly, we consider all rules that introduce new labels in a single phase, the \emph{label phase}.
We will show that the whole proof search is well defined by showing that each phase terminates. 
The important thing here is that we do not rely on blind proof search as we want to ensure that every rule that we can apply (without creating duplicates) has to be applied eventually. This will simplify our countermodel construction as we can read out the valuation and the relational structure directly from the propositional and relational atoms.

\begin{definition}[Saturation]
    Let $\R , \Gamma \Ra \Omega$ be a sequent that contains a labelled formula $x:A$. We call $x:A$ \emph{saturated} in $\R , \Gamma \Ra \Omega$ if the following conditions hold based on the form of $A$:
    \begin{itemize}
        \item[-] $A$ is a propositional atom $p \in \Prop$ or $A= \bot$;
        \item[-] $A= A_1 \to A_2$ and $x :A \in \Gamma$, then either $x:A_1 \in \Omega$ or $x: A_2 \in \Gamma$;
        \item[-] $A= A_1 \to A_2$ and $x :A \in \Omega$, then $x:A_1 \in \Gamma$ and $x: A_2 \in \Omega$;
        \item[-] $A = \Box A_1$ and $x :A \in \Gamma$, then for all $y$ with $xRy \in \R$, $y :A_1 \in \Gamma$;
        \item[-] $A = \triangle A_1$ and $x :A \in \Gamma$, then for all $y$ with $xSy \in \R$, $y :A_1 \in \Gamma$.
    \end{itemize}
    A sequent is saturated, if all its formulas are saturated and $\R$ is closed under $\transBi$.
\end{definition}

We call a tree of sequents (finite or infinite) formed by rules of $\labCS$ (possibly containing some leaves that are not closed by axioms) a \emph{derivation}.
Let $\tree$ be a derivation with some open leaves. The \emph{saturation phase} over $\tree$ is then defined by applying the rules $\to$L, $\to$R, $\Box$L, $\triangle$L, and $\transBi$ as long as possible to the open leaves which are not saturated.
It is not important in which order we apply these rules, as the following lemma shows.

\begin{lemma} \label{lem:sat}
    Any saturation phase yields a finite tree whose open leaves are saturated.
\end{lemma}
\begin{proof}
    By construction, the open leaves of a saturation phase will be saturated. Let $\Seq$ be some open leaf at the start of the saturation phase.
We show that the saturation phase terminates by showing that only finitely many formulas can be added during saturation. 
Observe that $\Seq $ only contains finitely many formulas and every new formula that gets added to $\Seq $ must be a subformula of another formula in $\Seq $. Therefore, for every label $x$ occurring in $\Seq $ there can only be finitely many new formulas $x:A$ that can be added. As the amount of labels is finite and does not increase, the total amount of formulas that are added in the saturation phase is finite. 
\end{proof}

Next, we define the second phase, the \emph{label phase}, of our proof search as simply applying the rules $\Box$R and $\triangle$R continually as long as possible. 
We call a sequent to which no such rules can be applied \emph{label saturated}.

\begin{lemma}
    Any label phase yields a finite tree whose leaves are label saturated.
\end{lemma}
\begin{proof}
    The lemma follows directly from the observation that sequents contain only finitely many formulas and any formula can only have a finite amount of nested $\Box$ and $\triangle$.
\end{proof}

Both lemmas show us that the following \emph{proof search strategy} is well defined:
Alternatingly execute saturation phases and label phases while applying Id and $\bot$ eagerly (i.e.~as soon as possible). 

We will use this proof search strategy to obtain a countermodel from an unprovable sequent. 
A sequent which is saturated and label saturated and to which neither Id nor $\bot$ can be upwards applied is called \emph{fully saturated}. Note that any open leaves which our proof search strategy creates must be fully saturated. Let us show in the following lemma how to extract a countermodel from such a sequent.

\begin{lemma} \label{lem:fulsat}
    If $\Seq$ is a fully saturated sequent, then there is a countermodel ${\Mw \nVdash \Seq}$. 
\end{lemma}
\begin{proof}
    Let us first consider a fully saturated sequent $\R , \Gamma \Ra \Omega$.
    Define a model $\Mw = \la W , \prec , M_0 , M_1 , V \ra$ as follows: 
    \begin{itemize}
        \item[] $W := Lab ( \R , \Gamma \Ra \Omega)$;
        \item[] $M_{0} := \{ x \in W \spa \text{there is some } y \text{ such that } yRx \in \R \}$;
        \item[] $M_{1}:= \{ x \in W \spa \text{there is some } y \text{ such that } ySx \in \R \}$;
        \item[] ${\prec} := \{ (x,y) \in W \times W \spa xRy \in \R \text{ or } xSy \in \R \}$;
        \item[] $V(p) := \{ x \in W \spa x : p \in \Gamma \}$ for any $p \in \Prop$.
    \end{itemize}
    First, observe that $\la W , \prec, V \ra$ is a $\GL$-model, as $\R$ constitutes a finite and transitive tree by closure under $\transBi$, hence $\Mw$ is also a Carlson model. 
    
    Claim: $\Mw$ satisfies all formulas from $\Gamma$ and falsifies all formulas from $\Omega$ under the identity interpretation  $I(x)=x$ . 
    Consider any $x:A$ occurring in the sequent. We prove the claim by simultaneous induction on the complexity of $A$. 
    
    $A \in \Prop$: If $x:A \in \Gamma $, the claim follows by definition. If $x:A \in \Omega$, it follows from the fact that the sequent is saturated, so not an identity sequent. 
    
    $A=\bot$: We know that $x:A \in \Gamma$ cannot happen because of full saturation, whereas if $x:A \in \Omega$ the formula is falsified by definition.
    
    $A= A_1 \to A_2$: If $x:A \in \Gamma$, then by full saturation either $x: A_1 \in \Omega$ or $x:A_2 \in \Gamma$. So, by induction hypothesis $\Mw , x \nVdash A_1$ or $\Mw , x \Vdash A_2$, thus in both cases $\Mw , x\Vdash A_1 \to A_2$. If $x:A \in \Omega$, then $A_1 \in \Gamma$ and $A_2 \in \Omega$. By induction hypothesis $\Mw , y \Vdash A_1$ and $\Mw ,y \nVdash A_2$, so $\Mw ,y \nVdash A_1 \to A_2$.
    
    $A = \Box A_1$: If $x:A \in \Gamma$, then $z : A_1 \in \Gamma$ for any $z$ with $xR z \in \R$. $\Mw, x \Vdash \Box A_1$ holds vacuously if no such $z$ exists. By induction hypothesis, $\Mw , z \Vdash A_1$ and by the definition of $\prec$ and the set $M_0$: $\Mw , x \Vdash \Box A_1$. The case for $x:A \in \Omega$ cannot occur due to label saturation. 
    
    The case of $A = \triangle A_1$ works similarly to the previous one. 
\end{proof} 

Now, we apply our proof search to show how to create a countermodel from a failed proof that is obtained by the previously introduced strategy. 

\begin{theorem}[Completeness of $\labCS$] \label{thm:complCS}
    For any sequent $\Seq$: If $\CS \Vdash \Seq$ then $\labCS$ $\vdash \Seq$. Furthermore, if a sequent is not provable in $\labCS$, there is a countermodel $\Mw \nVdash \Seq$.
\end{theorem}
\begin{proof}    
    We prove the contrapositive.
    Assume $\labCS \nvdash \Seq$. Perform the proof search described above on $\Seq$ to obtain a derivation $\tree$. Consider some branch $(\Seq_i)_{i < \omega}$ in $\tree$ such that it either ends with a fully saturated sequent, or it is infinite but not progressing. Such a branch must exist by $\labCS \nvdash \Seq$ as we otherwise would have shown the opposite. 
    Write $\Seq_i = \R_i , \Gamma_i \Ra \Omega_i$ and define a sequence of models $(\Mw_i)_{i < \omega}$ with $\Mw_i = \la W_i , M_{0 \, i}, M_{1 \, i}, \prec_i, V_i \ra$ as follows.
    \begin{itemize}
        \item[] $W_i := Lab(\Seq_i)$
        \item[] $M_{0 \, i} := \{ x \in W_i \spa \text{there is some } y \text{ such that } yRx \in \R_i \}$
        \item[] $M_{1 \, i}:= \{ x \in W_i \spa \text{there is some } y \text{ such that } ySx \in \R_i \}$
        \item[] ${\prec_i} := \{ (x,y) \in W_i \times W_i \spa xRy \in \R_i \text{ or } xSy \in \R_i \}$
        \item[] $V_i(p) := \{ x \in W_i \spa x : p \in \Gamma \}$
    \end{itemize}
    
    If $(\Seq_i)_{i < \omega}$ is finite and ends with a fully saturated sequent $\Seq_n$, by lemma \ref{lem:fulsat} $\Mw_n \nVdash \Seq_n$. Then, by lemma \ref{lem:FalsityCS} also $\Mw_n \nVdash \Seq$.
    
    If $(\Seq_n)_{n < \omega}$ is infinite, then any potential trace in $(\Seq_n)_{n < \omega}$ must be eventually constant. Consider the limit model $\Mw = \la W , \prec , M_0 , M_1 , V \ra$ with $W : =\bigcup_{i < \omega} W_i$, ${\prec} : =\bigcup_{i < \omega} \prec_i$, $M_0 : =\bigcup_{i < \omega} M_{0 \, i}$, $M_1 : =\bigcup_{i < \omega} M_{1 \, i}$, and $V(p) : =\bigcup_{i < \omega} V_i (p)$. 
    By every trace of $(\Seq_n)_{n < \omega}$ being eventually constant, $\prec$ must be conversely wellfounded, otherwise an infinite $\prec$-chain would allow to obtain a progressing trace. (We do not consider the case for $\prec$ having a cycle as it must form a tree if the root of the derivation is cycle free.) Therefore, $\Mw$ must be a Carlson model. We now show that for any $\Seq_i$ in the branch we have $\Mw \nVdash \Seq_i$. We do this similarly to the proof of lemma \ref{lem:fulsat} by induction over the complexity of formulas occurring in $\Seq_i$, for all $\Seq_i$ simultaneously. 
    
    Consider any sequent $\Seq_n$ in the branch and a formula $A$ such that $x:A$  occurs in $\Seq_n$. It then suffices to show that $\Mw, x \Vdash A$ if $x:A \in \Gamma_n$ and $\Mw , y \nVdash B$ if $x:A \in \Omega_n$.

    $A \in \Prop$: If $x:A \in \Gamma $, the claim follows by definition. If $x:A \in \Omega$, it follows from the fact that the branch did not close during the proof strategy, thus Id cannot be applied.

    $A=\bot$: We know that $x:A \in \Gamma$ cannot happen because then the rule $\bot$ could have been applied, making the branch finite in the proof search strategy. If $x:A \in \Omega$, the formula is falsified by definition.

    \sloppy
    $A= A_1 \to A_2$: By the strategy $A$ must be eventually principal. 
    If ${x:A \in \Gamma}$, there is a sequent $\Seq_m$ with ${n < m}$ such that either ${x: A_1 \in \Omega_m}$ or ${x:A_2 \in \Gamma_m}$.
    By induction hypothesis $\Mw , x \nVdash A_1$ or $\Mw , x \Vdash A_2$, thus in both cases ${\Mw , x\Vdash A_1 \to A_2}$. 
    If $x:A \in \Omega$, there is a sequent $\Seq_m$ with $n < m$ such that $A_1 \in \Gamma_m$ and $x: A_2 \in \Omega_m$. By induction hypothesis $\Mw , y \Vdash A_1$ and $\Mw ,y \nVdash A_2$, so $\Mw ,y \nVdash A_1 \to A_2$.
    
    $A = \Box A_1$: If $x:A \in \Gamma$, there must be some $\Seq_m$ with $n < m$ such that $z : A_1 \in \Gamma_m$ for any $z$ with $xR z \in \R_m$. $\Mw, x \Vdash \Box A_1$ holds vacuously if no such $z$ exists. By induction hypothesis, $\Mw , z \Vdash A_1$ and by the definition of $\prec$ and of the set $M_0$: $\Mw , x \Vdash \Box A_1$.
    If $x:A \in \Omega$, there is some $\Seq_m$ with $n < m$ such that $z : A_1 \in \Omega_m$ and $xRz \in \R_m$. By induction hypothesis $\Mw , z \nVdash A_1$, $y \prec z$, and $z \in M_0$. Therefore, $\Mw , y \nVdash \Box A_1$.
    
    The case of $A = \triangle A_1$ works similarly to the previous one. 

    As $\Mw \nVdash \Seq_i$ for all $i < \omega$, we gain as a corollary $\Mw \nVdash \Seq$.  
\end{proof}

This proof also yields directly a decision procedure for the logic $\CS$ as we can apply the strategy to a given sequent $\Ra x:A$ and either obtain a valid proof or extract a countermodel by the procedure described in the proof above.

We remark that we can make the infinite countermodel $\Mw$, which we obtained from a non-progressing branch, into a finite model:
observe that due to the model being infinite but conversely wellfounded we must have some points with infinitely many $\prec$-children; however, we consider only finitely many subformulas. Thus, if we quotient these $\prec$-children into equivalent points (meaning they satisfy the same formulas and are in the same subset of either $M_0$ or $M_1$), we can always obtain a finite set of $\prec$-children for every point without changing the valid formulas of the model. 

\begin{corollary}
    For any formula $A \in \LaBi$: $\CS \vdash A$ iff $\labCS \vdash A$.
\end{corollary}
\begin{proof}
    Apply theorems \ref{thm:soundCS} and \ref{thm:complCS} for the sequent $\Ra x: A$ and apply theorem \ref{thm:soundComplCS}.
\end{proof}

\section{Conclusion and Future Works}\label{sec:concl}
We have introduced a non-wellfounded calculus for the bimodal provability logic $\CS$ by a two-sorted labelled approach. We have shown soundness and completeness for this system, the latter of which relies on a primitive decision procedure and a countermodel construction. 

For the future, we want to investigate modal extensions of this calculus. This should be easy for extensions that have a first-order Kripke correspondence, as one can simply add the semantic conditions as rules for the relational atoms. 
For Kripke-incomplete logics, such as $\ER$ (see \cite{Carlson86,Visser95}), we conjecture that one could introduce a different progress condition which can capture the generalised semantics.

Building on that, it might be interesting to see whether these labelled proof systems might be reducible to simple sequent calculi without labels in a similar way to \cite{Lyon_2025,Pimentel19}, except that we would consider non-wellfounded proofs. We expect to find a system similar to the non-wellfounded one in \cite{Shamkanov14}.
For applications of these calculi, it will be interesting to see how they enable one to obtain results about these logics, such as interpolation (see \cite{Afsharietal_2021,Shamkanov14}).

Another line of further research could include the development of calculi for other provability logics such as Feferman provability \cite{Feferman60,Shavrukov94,Visser89}, slow provability \cite{Henk16}, intuitionistic provability \cite{Ardeshir18}, or polymodal provability logic $\GLP$ \cite{Japaridze85}.

\bibliographystyle{splncs04}
\bibliography{references}


\end{document}
